% =====================================================================================
% CALM-R3 paper block. Fill [[...]] from calm_r3/phase0_results.json (+ phase1_routed.json).
% Use ONLY if Phase 0 prints GO. If NO-GO, see the "negative-result" paragraph at the end.
% =====================================================================================

\subsection{Reflective routing: learning when not to reason}
Two of our findings point the same way: confidence predicts correctness, and a single
calibrated pass beats five-sample self-consistency. Together they suggest a policy rather
than a fixed pipeline---reason only when the fast answer is in doubt. We route on the
calibrated confidence of Section~\ref{sec:selection}. When the top option clears a
threshold $\tau$ the fast logit read is returned directly; otherwise the question is sent
to a reasoning path that produces a short differential before re-reading the option
logits. A second LoRA adapter, trained on the hard slice with a correction objective
\begin{equation}
\mathcal{L}=\mathrm{CE}(p_{\text{ref}},y)\;\mathbb{1}[\text{fast wrong}]
+\mu\,\mathrm{KL}\!\big(\overline{p_{\text{fast}}}\,\Vert\,p_{\text{ref}}\big)\,\mathbb{1}[\text{fast right}],
\end{equation}
supplies that path, while the original adapter serves the fast path; both are tiny, so the
pair fits one T4. The objective supervises reasoning toward the gold answer exactly where
the fast pass fails, and otherwise keeps it consistent with the confident fast prediction.

\subsection{Adaptive compute}
Because only a fraction of questions are routed to reasoning, the system trades accuracy
against decoding cost along a frontier. We sweep $\tau$ and report routed accuracy against
the mean number of generated tokens per question (Fig.~\ref{fig:r3}). At
$\tau=[[best_tau]]$, CALM-R3 reaches [[routed_acc]] accuracy while reasoning on only
[[frac_reasoned]]\% of items---[[delta_vs_fast]] points over the fast-only reader and
within [[gap_vs_reason]] of always-reasoning, at roughly [[token_pct]]\% of its decoding
tokens. The lesson for a small budget is that compute should follow uncertainty.

\begin{table}[tb]
\caption{Reflective routing on the 7B student (test subset, $n=[[n]]$). Fill from
\texttt{calm\_r3/phase0\_results.json} and \texttt{phase1\_routed.json}.}
\label{tab:r3}
\centering
\small
\begin{tabular}{@{}lccc@{}}
\toprule
Policy & acc & \% reasoned & tokens/q \\
\midrule
Always fast            & [[fast_acc]]   & 0   & 0 \\
Always reason          & [[reason_acc]] & 100 & [[reason_tok]] \\
\textbf{CALM-R3 ($\tau{=}[[best_tau]]$)} & \textbf{[[routed_acc]]} & [[frac_reasoned]] & [[routed_tok]] \\
\bottomrule
\end{tabular}
\end{table}

% --- contribution bullet ---
% \item A \textbf{confidence-routed reflective reader} that reasons only on uncertain
%  items, matching or beating the fast reader while spending a fraction of the decoding
%  tokens of an always-reason policy.

% --- IF PHASE 0 IS NO-GO, use this honest negative result instead of the above ---
% \subsection{When reasoning does not help}
% We tested a confidence-routed reflective reader that sends only low-confidence items to a
% chain-of-thought path. On this student the reasoning path did not exceed the fast read on
% the uncertain subset ([[reason_helps_bottom30]] absolute), so routing gave no accuracy
% gain. We read this as evidence that one-shot fine-tuning erodes the base model's
% reasoning, and that for this regime calibration and abstention---not extra reasoning---are
% the reliable levers.
